% ============================================================
%  PTE Core 模考 1B —— 写作 Writing 错题与详解·整理卷
%  版本 000   (self-contained / 单文件)
%  编译: xelatex 本文件.tex   (需 Noto CJK + TeX Gyre Heros 字体)
%  说明: 本文件由 build/paper/{ptestyle.sty,cover.tex,sec_writing_body.tex}
%        合并展平而成, 不依赖任何外部图片, 可独立编译复现同一 PDF。
% ============================================================
\documentclass[11pt]{article}

% ---------------- 样式 (原 ptestyle.sty, 已内联) ----------------
% ============================================================
%  ptestyle.sty  —  PTE 模考错题整理卷 共享样式
%  编译引擎: XeLaTeX  (中英混排, 需要 Noto CJK 字体)
% ============================================================

\usepackage[a4paper,margin=2cm,headsep=4mm]{geometry}
\usepackage{fontspec}
\usepackage{xeCJK}
\usepackage[table]{xcolor}
\usepackage[normalem]{ulem}        % \sout 删除线
\usepackage{tabularx}
\usepackage{array}
\usepackage{ragged2e}
\usepackage{enumitem}
\usepackage[most]{tcolorbox}
\usepackage{fancyhdr}
\usepackage{graphicx}
\usepackage{needspace}
\usepackage{tikz}

% ---------- 字体 ----------
\setmainfont{TeX Gyre Heros}            % 拉丁: Helvetica 风格(纤细清晰)
\setCJKmainfont{Noto Sans CJK SC}
\setCJKsansfont{Noto Sans CJK SC}
\newfontfamily\cjkserif{Noto Serif CJK SC}
\linespread{1.18}
\setlength{\parindent}{0pt}
\setlength{\parskip}{4pt}

% ---------- 配色 ----------
\definecolor{cWrong}{HTML}{D32F2F}     % 红  = 修改 / 错误
\definecolor{cFix}{HTML}{1B7F3B}       % 绿  = 正确建议
\definecolor{cDelete}{HTML}{1565C0}    % 蓝  = 删除
\definecolor{cInsert}{HTML}{7B1FA2}    % 紫  = 插入
\definecolor{cPartial}{HTML}{E07B00}   % 橙  = 部分得分
\definecolor{cNote}{HTML}{0277BD}      % 蓝  = 注解
\definecolor{cAccent}{HTML}{16A085}    % 主题青绿(平台主色)
\definecolor{cWrongBg}{HTML}{FBE0E0}
\definecolor{cDelBg}{HTML}{DCE8FB}
\definecolor{cInsBg}{HTML}{EEE0F5}
\definecolor{cHead}{HTML}{20303A}
\definecolor{cGrayBox}{HTML}{F4F6F7}
\definecolor{cModelBg}{HTML}{EAF7EF}
\definecolor{cYouBg}{HTML}{FFF7F5}
\definecolor{cNoteBg}{HTML}{EAF3FA}
\definecolor{cRule}{HTML}{D7DEE2}

% ---------- 行内批改宏 (复刻平台标注) ----------
% \fix{原词}{正确}  红色删除线 + 绿色上标改正  (修改/拼写/空格)
\newcommand{\fix}[2]{%
  \colorbox{cWrongBg}{\textcolor{cWrong}{\sout{#1}}}%
  \,\textsuperscript{\textcolor{cFix}{\footnotesize #2}}}
% \del{原词}  蓝色删除线  (应删除)
\newcommand{\del}[1]{%
  \colorbox{cDelBg}{\textcolor{cDelete}{\sout{#1}}}%
  \,\textsuperscript{\textcolor{cDelete}{\footnotesize 删}}}
% \ins{词}  紫色下划线  (应插入)
\newcommand{\ins}[1]{%
  \textsuperscript{\textcolor{cInsert}{\footnotesize $\wedge$}}%
  \colorbox{cInsBg}{\textcolor{cInsert}{#1}}}
% \miss{词}  灰色删除线  (口语漏读 / 未作答)
\newcommand{\miss}[1]{\textcolor{gray}{\sout{#1}}}
% \add{原词}{改正}  橙色波浪线 + "补"标签  (平台漏标、本卷补标的错误)
\definecolor{cAdd}{HTML}{E65100}
\newcommand{\add}[2]{%
  \textcolor{cAdd}{\uwave{#1}}%
  \,\textsuperscript{\textcolor{cAdd}{\footnotesize 补：#2}}}
% 高亮(口语发音): 好/中/差
\newcommand{\spk}[2]{\textcolor{#1}{#2}}

% ---------- 评分单元格着色 ----------
\newcommand{\sfull}[1]{\textcolor{cFix}{\bfseries #1}}     % 满分
\newcommand{\spart}[1]{\textcolor{cPartial}{\bfseries #1}} % 部分
\newcommand{\szero}[1]{\textcolor{cWrong}{\bfseries #1}}   % 零/低分

% ---------- 分数小标签 ----------
\newtcbox{\chip}[1][cAccent]{on line, boxsep=2pt, left=4pt, right=4pt, top=1pt,
  bottom=1pt, arc=3pt, colback=#1, colframe=#1, coltext=white, nobeforeafter,
  fontupper=\bfseries\footnotesize}

% ---------- 题块小标题 ----------
\newcommand{\ptesub}[1]{%
  \par\needspace{2\baselineskip}\smallskip
  {\color{cAccent}\rule[-0.2ex]{3pt}{1.05em}}\hspace{6pt}%
  {\bfseries\color{cHead}#1}\par\nopagebreak\smallskip}

% ---------- 题块外框 ----------
% \begin{qblock}{标号·题型}{右上信息(得分/用时)} ... \end{qblock}
\newtcolorbox{qblock}[2]{
  breakable, enhanced, arc=2.5pt,
  colback=white, colframe=cRule, boxrule=0.6pt,
  left=11pt, right=11pt, top=8pt, bottom=10pt,
  toptitle=3pt, bottomtitle=3pt, lefttitle=11pt, righttitle=11pt,
  colbacktitle=cHead, coltitle=white, fonttitle=\bfseries,
  title={#1\hfill\normalfont\bfseries\small #2}}

% ---------- 文本块: 题干 / 范文 / 你的作答 / 建议 ----------
\newtcolorbox{promptbox}{enhanced, breakable, colback=cGrayBox, colframe=cGrayBox,
  arc=2pt, left=8pt, right=8pt, top=6pt, bottom=6pt, boxrule=0pt}
\newtcolorbox{modelbox}{enhanced, breakable, colback=cModelBg, colframe=cModelBg,
  borderline west={3pt}{0pt}{cFix}, arc=1pt, left=8pt, right=8pt, top=6pt, bottom=6pt, boxrule=0pt}
\newtcolorbox{youbox}{enhanced, breakable, colback=cYouBg, colframe=cYouBg,
  borderline west={3pt}{0pt}{cWrong}, arc=1pt, left=8pt, right=8pt, top=6pt, bottom=6pt, boxrule=0pt}
\newtcolorbox{notebox}{enhanced, breakable, colback=cNoteBg, colframe=cNoteBg,
  borderline west={3pt}{0pt}{cNote}, arc=1pt, left=8pt, right=8pt, top=6pt, bottom=6pt, boxrule=0pt}

% ---------- 评分表 ----------
% 用法见正文; 这里给表头宏
\newcommand{\scorehead}{%
  \rowcolors{2}{cGrayBox}{white}
  \arrayrulecolor{cRule}}

% ---------- 章节大标题 ----------
\newcommand{\ptesection}[3]{% #1=中文名 #2=英文名 #3=该项分数
  \needspace{6\baselineskip}
  \begin{tcolorbox}[enhanced, colback=cAccent, colframe=cAccent, sharp corners,
     left=14pt, right=14pt, top=10pt, bottom=10pt, boxrule=0pt]
    {\fontsize{20}{24}\selectfont\bfseries\color{white} #1\ \ }%
    {\large\color{white!85} #2}\hfill
    {\large\color{white}本项得分\ \ }{\fontsize{22}{24}\selectfont\bfseries\color{white} #3}%
    \ {\color{white!85}/ 90}
  \end{tcolorbox}}

% ---------- 颜色图例 ----------
\newcommand{\ptelegend}{%
  \begin{tcolorbox}[enhanced, colback=cGrayBox, colframe=cRule, boxrule=0.5pt,
    arc=2pt, left=10pt, right=10pt, top=6pt, bottom=6pt]
  {\bfseries\color{cHead}颜色说明\quad}
  \fix{改}{正}~修改/拼写\quad
  \del{删}~应删除\quad
  \ins{插}~应插入\quad
  \add{补}{改正}~平台漏标(本卷补标)\quad
  {\color{cFix}\rule{1em}{0.6em}}~范文/正确\quad
  {\color{cNote}\rule{1em}{0.6em}}~AI 建议
  \end{tcolorbox}}

% ---------- 页眉页脚 ----------
\pagestyle{fancy}
\fancyhf{}
\renewcommand{\headrulewidth}{0.4pt}
\renewcommand{\footrulewidth}{0pt}
\fancyhead[L]{\footnotesize\color{cHead}\bfseries PTE Core 模考 1B}
\fancyhead[R]{\footnotesize\color{gray} 错题与详解整理卷}
\fancyfoot[C]{\footnotesize\color{gray}\thepage}

\begin{document}

% ---------------- 封面 (原 cover.tex) ----------------
% ---------- 封面 / 成绩总览 ----------
\definecolor{mListen}{HTML}{1F3A5F}
\definecolor{mRead}{HTML}{C2D70B}
\definecolor{mSpeak}{HTML}{4D4D4D}
\definecolor{mWrite}{HTML}{A01B7D}

\begin{titlepage}
\thispagestyle{empty}
\centering
\vspace*{1.4cm}

{\fontsize{30}{36}\selectfont\bfseries\color{cHead} PTE Core 模考 1B}\\[4mm]
{\fontsize{17}{20}\selectfont\color{cAccent}\bfseries 错题与详解 · 整理卷}\\[3mm]
{\color{gray} APEUni / 猩际 PTE 模考　·　提交时间 2026-06-21 12:36}\\[1.2cm]

% ---- 总分 ----
\begin{tcolorbox}[width=0.62\textwidth, halign=center, enhanced,
  colback=cAccent, colframe=cAccent, arc=6pt, boxrule=0pt, top=10pt, bottom=10pt]
  {\color{white}\large 总分 Overall}\\[2pt]
  {\fontsize{46}{50}\selectfont\bfseries\color{white} 40}
\end{tcolorbox}

\vspace{1.0cm}

% ---- 四项分数条形图 ----
\begin{tcolorbox}[width=0.84\textwidth, enhanced, colback=white, colframe=cRule,
  boxrule=0.8pt, arc=4pt, left=18pt, right=18pt, top=14pt, bottom=14pt]
\setlength{\tabcolsep}{6pt}\renewcommand{\arraystretch}{1.7}
\sffamily
\begin{tabular}{@{}r@{\hskip 8pt}l@{\hskip 10pt}l@{}}
 \bfseries Listening 听力 & {\color{mListen}\rule{3.67cm}{0.42cm}} & {\bfseries\large 33}\\
 \bfseries Reading 阅读   & {\color{mRead}\rule{4.67cm}{0.42cm}}   & {\bfseries\large 42}\\
 \bfseries Speaking 口语  & {\color{mSpeak}\rule{3.22cm}{0.42cm}}  & {\bfseries\large 29}\\
 \bfseries Writing 写作   & {\color{mWrite}\rule{6.33cm}{0.42cm}}  & {\bfseries\large 57}\\
\end{tabular}\\[2mm]
{\footnotesize\color{gray} 满分 90　|　刻度: 条长 = 得分 / 90}
\end{tcolorbox}

\vfill
\begin{flushleft}\small\color{gray}
\textbf{说明:}本卷由本次模考的题目与"详解"截图逐题誊写、重新排版而成,完整保留每道题的
\emph{题干 / 原文、范文、你的作答、逐项评分、彩色批改与 AI 中文建议}。\\
所有错误按平台标注用颜色标出(见下方图例)。原始"详解"PDF 不可复制、仅截图,本卷内容均为人工对照誊录,若与原图有出入以原图为准。
\end{flushleft}
\vspace{4mm}
\ptelegend
\end{titlepage}

% ---------------- 写作正文 (原 sec_writing_body.tex) ----------------
% ===================== 写作 Writing =====================
\ptesection{写作}{Writing · Summarize Written Text + Write Email}{57}

\vspace{1mm}
{\small\color{gray} 本部分含 \textbf{SWT 概括写作} 2 题（满分各 8）与 \textbf{WE 邮件写作} 2 题（满分各 15）。
SWT 仅展示“你的作答 + 批改”；WE 另含平台“范文”。下方每题保留逐项评分与 AI 中文建议。}
\vspace{2mm}

\newcolumntype{Y}{>{\RaggedRight\arraybackslash}X}

% ---------------- Q1 ----------------
\begin{qblock}{Q1\quad SWT \#57\ \ Metaverse}{7.5 / 8\quad·\quad 用时 09:06}

\ptesub{原文 Source Text}
\begin{promptbox}
The concept of the metaverse is rapidly gaining traction in the world of technology, promising to revolutionize the way we interact with digital spaces and each other. It represents a new era of connectivity and immersion, where physical and digital boundaries blur. In essence, the metaverse is a virtual, interconnected universe where users can engage in various activities, interact with others, and create their own digital presence. It encompasses virtual reality (VR), augmented reality (AR), and the internet, seamlessly blending these technologies into a unified experience.

One of the most exciting aspects of the metaverse is its potential to transform the way we work, socialize, learn, and entertain ourselves. Imagine attending a virtual meeting in a lifelike conference room, exploring distant places in an immersive travel experience, or attending a concert with friends from around the world --- all from the comfort of your home. Businesses are already exploring the possibilities of the metaverse, from virtual showrooms for products to interactive training simulations for employees. It has the potential to disrupt traditional industries and create new opportunities for innovation and entrepreneurship.
\end{promptbox}

\ptesub{你的作答 Your Response}
\begin{youbox}
The concept of the metaverse is gaining traction in the world of technology, and it represents a new era of connectivity and immersion, and the metaverse has potential to \fix{transfor}{transform} the way we work, and it can disrupt traditional industries and create opportunities for innovation and entrepreneurship.
\end{youbox}

\ptesub{评分明细 Score Breakdown}
{\small\scorehead
\begin{tabularx}{\linewidth}{l c Y}
\rowcolor{cHead}\textcolor{white}{\bfseries 维度} & \textcolor{white}{\bfseries 得分} & \textcolor{white}{\bfseries AI 建议}\\
Content 内容        & \sfull{2 / 2}   & 内容很棒\\
Form 格式          & \sfull{2 / 2}   & Form 很棒，这是必拿分数\\
Grammar 语法       & \sfull{2 / 2}   & 很棒，AI 小猩几乎检查不出什么语法问题\\
Vocabulary 词汇    & \spart{1.5 / 2} & SWT(C) 中可以摘抄原文，但记得摘抄时也要注意拼写以及细节语法点哦\\
\end{tabularx}}
\smallskip
\textbf{本题得分：}\ {\color{cAccent}\bfseries 7.5 / 8}\quad{\footnotesize\color{gray}（错误：\fix{transfor}{transform} —— 拼写错误）}

\end{qblock}

% ---------------- Q2 ----------------
\begin{qblock}{Q2\quad SWT \#56\ \ Transformative Power}{7.3 / 8}

\ptesub{原文 Source Text}
\begin{promptbox}
In today's world, cinema and television have become an integral part of our lives, impacting us in various profound ways. From entertainment and education to cultural expression, they hold immense significance.

Cinema and television offer a window into different cultures and realities. Through films and TV shows, we can explore diverse perspectives, lifestyles, and experiences. This exposure fosters empathy and a broader worldview, helping break down stereotypes and prejudices.

Furthermore, these media platforms have the power to educate. Documentary films and educational programs provide valuable insights into historical events, scientific discoveries, and social issues. They make complex subjects accessible and engage audiences of all ages.

Cinema and television are also powerful tools for cultural expression. Filmmakers and TV creators use their work to address societal concerns, challenge norms, and promote change. Iconic movies like ``To Kill a Mockingbird'' and TV series like ``The Handmaid's Tale'' have sparked crucial conversations about injustice and inequality.
\end{promptbox}

\ptesub{你的作答 Your Response}
\begin{youbox}
Cinema and television have become an integral part of our lives, and \fix{the}{they} offer a window into different cultures and realities, and they are also powerful tools for cultural expression, and they have the power to educate.
\end{youbox}

\ptesub{评分明细 Score Breakdown}
{\small\scorehead
\begin{tabularx}{\linewidth}{l c Y}
\rowcolor{cHead}\textcolor{white}{\bfseries 维度} & \textcolor{white}{\bfseries 得分} & \textcolor{white}{\bfseries AI 建议}\\
Content 内容        & \sfull{2 / 2}    & 内容很棒\\
Form 格式          & \sfull{2 / 2}    & Form 很棒，这是必拿分数\\
Grammar 语法       & \spart{1.5 / 2}  & 点击有颜色标注的文字查看语法错误细节哦~\\
Vocabulary 词汇    & \spart{1.75 / 2} & SWT(C) 中可以摘抄原文，但记得摘抄时也要注意拼写以及细节语法点哦\\
\end{tabularx}}
\smallskip
\textbf{本题得分：}\ {\color{cAccent}\bfseries 7.25 / 8}\ {\footnotesize\color{gray}（平台显示 7.3）}\quad{\footnotesize\color{gray}（错误：\fix{the}{they}）}

\end{qblock}

% ---------------- Q3 ----------------
\begin{qblock}{Q3\quad WE \#20\ \ Cultivating Community Spirit through Engagement}{10.8 / 15\quad·\quad 用时 08:52}

\ptesub{题干 Task Prompt}
\begin{promptbox}
Your neighborhood association is looking to enhance local engagement and foster a sense of community among residents. Write an email to the association president suggesting three activities that could bring neighbors together and strengthen community bonds. You should write between 80 and 120 words.

\smallskip Your suggestions should focus on the following three themes:
\begin{itemize}[leftmargin=1.4em,topsep=2pt,itemsep=0pt]
  \item A neighborhood clean-up day;
  \item A community potluck or food fair;
  \item Regular sports or game nights.
\end{itemize}
You should include all three themes. Provide supporting ideas for your suggestions.
\end{promptbox}

\ptesub{范文 Model Answer（平台参考）}
\begin{modelbox}
Dear Association President,

I hope this message finds you well. In light of our shared goal to foster a stronger sense of community within our neighborhood, I propose three activities aimed at bringing residents closer together.

First, organizing a neighborhood clean-up day could not only beautify our area but also instill a collective sense of pride and accomplishment.

Second, hosting a community potluck or food fair would provide a wonderful opportunity for neighbors to share meals and stories, promoting cultural exchange and understanding.

Lastly, setting up regular sports or game nights can offer a fun and relaxed setting for people of all ages to interact and build friendships.

Warm regards,\\ Peter Pan
\end{modelbox}

\ptesub{你的作答 Your Response（含批改）}
\begin{youbox}
Dear association president,

\fix{I'm pleasure}{I'm happy} to write \fix{a}{an} email to you. I want to share my suggestions about enhancing local engagement and \fix{fosterring}{fostering} a \fix{sence}{sense} of community among residents.

First of all, I suggest that we can have a neighborhood clean-up day so that the environment will be \del{more} clean.

Then, I suggest that we can have a community potluck or food fair so that the community will be more peaceful.

Last but not least, I suggest that we can have \fix{regularsports}{regular sports} or game night so that the community will be more peaceful.

That's all. Looking forward to your \add{replay}{reply}.

Best \fix{recards}{regards},\\ Kathy
\end{youbox}
{\footnotesize\color{gray} 注：橙色波浪线（补）为平台\textbf{漏标}、本卷补标的错误。``replay'' 应为 ``reply''。}

\ptesub{评分明细 Score Breakdown}
{\small\scorehead
\begin{tabularx}{\linewidth}{l c Y}
\rowcolor{cHead}\textcolor{white}{\bfseries 维度} & \textcolor{white}{\bfseries 得分} & \textcolor{white}{\bfseries AI 建议}\\
Content 内容              & \sfull{3 / 3}   & 内容很棒\\
Form 格式                & \sfull{2 / 2}   & Form 很棒，这是必拿分数\\
Email Conventions 邮件格式 & \sfull{2 / 2}   & 邮件书写格式很棒！\\
Organization 结构         & \sfull{2 / 2}   & 段落分明，文章架构清晰明了，很棒\\
Vocabulary 词汇           & \spart{1 / 2}   & PTE 对词汇不要求很难，但需要恰当且拼写正确。点击彩色单词查看错误和修改建议\\
Grammar 语法              & \spart{0.8 / 2} & 语法在 PTE 写作中占分比很重，但其实又非常容易提高。丢分这么严重，有点不应该哦。建议找老师咨询下迅速提高的技巧\\
Spelling 拼写             & \szero{0 / 2}   & 你犯了 5 个拼写错误。太多了，得细心点啊\\
\end{tabularx}}
\smallskip
\textbf{本题得分：}\ {\color{cAccent}\bfseries 10.8 / 15}\quad{\footnotesize\color{gray}（平台标注：I'm pleasure→I'm happy、a→an、fosterring/sence/recards 拼写、regularsports 缺空格、more 应删除；\textcolor{cAdd}{补标：replay→reply}）}

\end{qblock}

% ---------------- Q4 ----------------
\begin{qblock}{Q4\quad WE \#21\ \ Revitalizing Library Engagement}{12.2 / 15\quad·\quad 用时 09:00}

\ptesub{题干 Task Prompt}
\begin{promptbox}
Your city's public library system is seeking ways to boost library membership and encourage more residents to take advantage of its resources. As a frequent visitor and library enthusiast, you want to offer suggestions to the library's management. Write an email to Ms.\ Bell, the director of library services, proposing three ideas to increase library engagement among city residents. You should write between 80 and 120 words.

\smallskip Your suggestions should focus on the following three themes:
\begin{itemize}[leftmargin=1.4em,topsep=2pt,itemsep=0pt]
  \item An annual book festival;
  \item Collaborations with local authors;
  \item Innovative membership benefits.
\end{itemize}
You should include all three themes. Provide supporting ideas for your suggestions.
\end{promptbox}

\ptesub{范文 Model Answer（平台参考）}
\begin{modelbox}
Dear Ms.\ Bell,

I'm excited to share ideas to enhance library engagement across our community.

{\color{cWrong}\footnotesize[此处原始截图在分页处缺失一句（Firstly… 关于 annual book festival 的建议），原图未截全。]}

Secondly, collaborations with local authors for workshops or storytelling sessions could offer unique insights into the writing process, encouraging budding writers and readers alike.

Lastly, introducing innovative membership benefits, such as exclusive online resources or early access to new arrivals, might incentivize more residents to join and actively use the library.

I believe these initiatives could significantly boost library membership and community involvement.

Warm regards,\\ Peter Pan
\end{modelbox}

\ptesub{你的作答 Your Response（含批改）}
\begin{youbox}
Dear Ms.\ Bell,

\fix{I'm pleasure}{I'm happy} to write an email to you. I want to offer some suggestions about \del{the} ways to boost library membership and encourage more residents to take advantage of its resources.

Firstly, my suggestion is that we can hold an annual book festival so that we can encourage more residents to take advantage of its resources.

Secondly, my suggestion is that we can build collaborations with local authors so that we can encourage more authors to participate.

Lastly, my suggestion is that we can \add{innovative}{前缺动词，如 introduce} membership benefits so that we can boost library membership.

That's all. \fix{I'mlooking}{I'm looking（缺空格）} forward to your reply.

Best regards,\\ Kathy
\end{youbox}

\ptesub{评分明细 Score Breakdown}
{\small\scorehead
\begin{tabularx}{\linewidth}{l c Y}
\rowcolor{cHead}\textcolor{white}{\bfseries 维度} & \textcolor{white}{\bfseries 得分} & \textcolor{white}{\bfseries AI 建议}\\
Content 内容              & \sfull{3 / 3}   & 内容很棒\\
Form 格式                & \sfull{2 / 2}   & Form 很棒，这是必拿分数\\
Email Conventions 邮件格式 & \sfull{2 / 2}   & 邮件书写格式很棒！\\
Organization 结构         & \sfull{2 / 2}   & 段落分明，文章架构清晰明了，很棒\\
Vocabulary 词汇           & \spart{1 / 2}   & PTE 对词汇不要求很难，但需要恰当且拼写正确。点击彩色单词查看错误和修改建议\\
Grammar 语法              & \spart{1.2 / 2} & 点击有颜色标注的文字查看语法错误细节哦~\\
Spelling 拼写             & \spart{1 / 2}   & 你犯了 2 个拼写错误\\
\end{tabularx}}
\smallskip
\textbf{本题得分：}\ {\color{cAccent}\bfseries 12.2 / 15}\quad{\footnotesize\color{gray}（平台标注：I'm pleasure→I'm happy、the 应删除、I'mlooking 缺空格；\textcolor{cAdd}{补标：innovative 前缺动词 introduce}）}

\end{qblock}

\end{document}
